\chapter{Experimental Setup}
\label{ch:ExperimentalSetup}

This chapter describes the datasets used for the analyses conducted in \cref{ch:metricsAnalysis}
and the software tools developed to create these datasets.

\section{Datasets}
\label{sec:datasets}

To carry out the analyses conducted in \cref{ch:metricsAnalysis}, we created two datasets:
one Goodware (\cref{sec:goodwareDataset}) and one Malware (\cref{sec:malwareDataset}).

These two datasets are a collection of metrics from a set of \ac{NPM} packages.

The metrics considered, discussed in the relevant section in \cref{ch:metricsAnalysis}, are identical for both datasets.
To compute this metric from the \ac{NPM} packages, we developed a software tool described in \cref{sec:ToolForAnalyzingNpmPackages}.

Instead, the set of \ac{NPM} packages and their collection are different for the two datasets.
However, both datasets always have 20 versions for each package. 
This was done to study their behavior and evolution across versions.

Consequently, these datasets consist of a collection of CSV files, one for each metric, 
containing the values of that specific metric for every version of every package included in the dataset.
We used the tool described in \cref{sec:ToolForAnalyzingNpmPackages} to generate both 
datasets, available on GitHub\footnote{\url{https://github.com/Frazzerz/Datasets-NPM-Analysis}}.

The two datasets are comparable to each other 
as they are homogeneous, since both have packages with the same number of versions (20).

\subsection{Goodware dataset}
\label{sec:goodwareDataset}


This dataset is composed of packages ranked among the 1,000 most popular 
in the \ac{NPM} registry, in terms of number of stars~\cite{tristan2026listtop10000npmpackages}.

Out of these 1,000 packages, 371 have less than 20 versions 
available in the \ac{NPM} registry.
Therefore, the dataset includes the 
remaining 629 packages, with their 20 most recent available versions parsable via node-semver~\cite{semver}.

To automatically obtain these versions of these packages, we used the developed tool (\cref{sec:ToolForAnalyzingNpmPackages}),
specifying the first 1,000 packages names from list of \textcite{tristan2026listtop10000npmpackages} in a JSON file as input.

\subsection{Malware dataset}
\label{sec:malwareDataset}


This dataset is composed of 20 \ac{NPM} packages
compromised by one of the four \ac{SSCA} analyzed in 
\cref{ch:AnalysisRecentAttacks}.

Each of these packages has 20 versions: one malicious and 19 legitimate.

To obtain the malicious versions of these packages, we used the developed script 
explained in \cref{sec:DownloadMaliciousTarballs}.

The list below reports the compromised packages in the Malware dataset for each attack,
including the specific malicious version number.
\begin{itemize}
    \item For the attack A (\cref{sec:attackA}) there are 3 packages:
    \begin{itemize}
        \item \texttt{eslint-config-prettier}, version 10.1.7
        \item \texttt{eslint-plugin-prettier}, version 4.2.3
        \item \texttt{synckit}, version 0.11.9
    \end{itemize}
    \item For the attack B (\cref{sec:attackB}) there are 9 packages:
    \begin{itemize}
        \item \texttt{ansi-regex}, version 6.2.1
        \item \texttt{ansi-styles}, version 6.2.2
        \item \texttt{chalk}, version 5.6.1
        \item \texttt{color-convert}, version 3.1.1
        \item \texttt{color-string}, version 2.1.1
        \item \texttt{debug}, version 4.4.2
        \item \texttt{supports-color}, version 10.2.1
        \item \texttt{strip-ansi}, version 7.1.1
        \item \texttt{wrap-ansi}, version 9.0.1
    \end{itemize}
    \item For the attack C (\cref{sec:attackC}) there are 4 packages:
        \begin{itemize}
            \item \texttt{@ctrl/deluge}, version 7.2.1
            \item \texttt{@ctrl/shared-torrent}, version 6.3.2
            \item \texttt{@ctrl/tinycolor}, version 4.1.2
            \item \texttt{ember-browser-services}, version 5.0.3
        \end{itemize}
    \item For the attack D (\cref{sec:attackD}) there are 4 packages:
        \begin{itemize}
            \item \texttt{@asyncapi/avro-schema-parser}, version 3.0.26
            \item \texttt{@asyncapi/bundler}, version 0.6.6
            \item \texttt{@asyncapi/cli}, version 4.1.3
            \item \texttt{posthog-node}, version 5.11.3
        \end{itemize}
\end{itemize}

Attacks C and D involved a large number of compromised packages, respectively 187 and 1,055.
For this reason, it was decided to include in the dataset a limited number of 
the most relevant packages for both attacks,
specifically those with a high number of weekly downloads (at least over 1,000).

The Malware dataset only includes packages with at least 19 legitimate versions available in the \ac{NPM} registry.

However in these attacks, there are compromised packages that have fewer than 19 versions. 
In such case, they were ignored and not included in the dataset (nor are they present in the list above).
Specifically, two packages were excluded from attack A and nine from attack B.

For the purpose of this dataset and its analysis conducted in \cref{ch:metricsAnalysis},
it is not necessary that its packages are among the 1,000 most popular in the \ac{NPM} registry.
Anyway in the list, all nine packages from attack B are among the 1,000 most popular,
while for the other attacks, none of the compromised packages fall into that category.